problem:  
Let \( (X^n, d) \) be a compact Alexandrov space of curvature \(\ge 0\) carrying an isometric, effective action by a compact Lie group \(G\).  
Suppose the regular orbit space \(Y = X/G\) is a smooth Riemannian manifold of nonnegative sectional curvature, and the projection \(\pi : X \to Y\) is a *submetry* whose fibers are principal \(G\)-orbits.  
Let \(\Sigma_x X\) denote the space of directions at \(x\), and let \(\rho_x : G_x \curvearrowright \Sigma_x X\) be the induced isotropy representation at \(x\).

Assume that for every point \(x \in X\),  
1. \( \Sigma_x X^{G_x} \) (the fixed-point set of the isotropy action on the space of directions) has constant intrinsic curvature \(1\), and  
2. the variation of these isotropy data \(x \mapsto (\Sigma_x X, \rho_x)\) is Hölder continuous in the pointed Gromov–Hausdorff topology.  

**Question:**  
Does such controlled isotropy-variation imply that \(X\) arises as the *equivariant* Gromov–Hausdorff limit of smooth, nonnegatively curved Riemannian \(G\)-manifolds with corresponding isotropy types?  
Equivalently, do Hölder-continuous Alexandrov isotropy cones force *equivariant smoothability* under nonnegative curvature?  

---

Why is it a "good" problem:  

1. **Conceptual depth:**  
   This question goes beyond the existing slice and recognition theorems in Alexandrov geometry. Those theorems address the *local* topology of a fixed \(G\)-orbit but do not control how isotropy representations change along the orbit space. Here, the focus shifts to the *quantitative regularity* (Hölder continuity) of isotropy variation and its potential to enforce global smooth realizability. The problem thus proposes a genuinely new analytic-to-topological transition principle in equivariant Alexandrov geometry.

2. **Bridge between fields:**  
   It connects:  
   - **Alexandrov geometry:** via curvature \(\ge0\) and tangent cones with group actions;  
   - **Geometric analysis:** through Hölder control on isotropy metrics and the cone structure;  
   - **Transformation group theory:** by requiring fiberwise smooth approximation under group symmetry;  
   - **Metric limit theory:** through equivariant Gromov–Hausdorff convergence.  
   Solving it would integrate metric comparison geometry, Lie group actions, and quantitative limit analysis into a unified framework.

3. **Potential significance:**  
   - A *positive answer* would yield a new **equivariant smoothing theorem**, unveiling a deep structural stability of nonnegatively curved Alexandrov G-spaces under quantitative isotropy control.  
   - A *negative answer*—a counterexample—would reveal new kinds of “quantitatively smooth but intrinsically singular” Alexandrov G-models, leading to previously unseen geometric phenomena under group symmetry.  
   Either outcome would redefine the boundary between smooth and singular symmetry in nonnegative curvature geometry, extending Perelman’s and Petrunin’s stability notions into the equivariant realm.

4. **Simplicity and profundity:**  
   The statement is crisp: “Do Hölder‑continuous isotropy cones guarantee equivariant smoothability?” Despite its short formulation, it encodes a high‑dimensional interaction among curvature, symmetry, and metric regularity. Its resolution would require new tools intertwining group actions, metric geometry, and quantitative analysis—an archetype of the “narrow entrance, profound depth” characteristic of a truly great mathematical problem.